\documentclass[10pt]{article}
\usepackage[margin=0.85in]{geometry}
\usepackage{amsmath,amssymb,amsthm}
\usepackage[colorlinks,urlcolor=blue]{hyperref}
\newtheorem*{problem1}{Problem 1}
\newtheorem*{problem2}{Problem 2}
\setlength{\parskip}{4pt}
\title{Two problems on rectangle blackouts}
\author{Arjun Pemmasani}
\date{15 November 2026}
\begin{document}
\maketitle

\section*{Setting}

Take the $h\times m$ grid of lattice points. A \emph{blackout} is a set $S$ of grid points such
that every rectangle with its 4 corners on grid points, tilted rectangles included, is determined
by its corners outside $S$. Let $f(h,m)$ be the largest size of a blackout. This is Problem~001
of Peter Kagey's Open Problem Collection (OEIS A397315).

Put $q_h=\lfloor h^2/4\rfloor$. For $m\ge3$ we have $f(h,m)\le m+q_h$. Let $m_h$ be the least
width such that $f(h,m)=m+q_h$ for every $m\ge m_h$. I can prove
\[
 q_h\le m_h\le C\,h^2\,\tau^*(h)^2,
\]
where $\tau^*(h)$ is the largest number of divisors of an integer up to $h^2$. I also have a draft
proof of $m_h\le C_\varepsilon h^2(\log h)^{1.3108+\varepsilon}$. The values are $m_4=6$, $m_5=8$,
$m_6=15$ and $m_7\le21$. Two questions remain. Is $m_h=O(h^2)$? Is $m_h>q_h$ for every $h$?

\section*{Problem 1. Bounded choices on a Hilbert curve}

My constructions place one pair of black points for each edge of a complete bipartite graph on
the rows. The main obstruction is a pair of points in rows of opposite parity at an integer
distance. The cleanest form of the question is the following.

Let $s=2^k$ and let $r(i,j)\in\{0,\dots,s^2-1\}$ be the Hilbert rank of the cell $(i,j)$,
$0\le i,j<s$. Fix integers $S\ge1$ and $M\ge1$. A \emph{choice} is a function
$\xi$ from the cells to $\{0,\dots,M-1\}$. Put $c(i,j)=M\,r(i,j)+\xi(i,j)$. An ordered pair of distinct cells
$u=(i_u,j_u)$ and $v=(i_v,j_v)$ \emph{collides} if
\[
 \bigl(S\,(c(v)-c(u))\bigr)^2+\bigl(2(j_v-i_u)-4s+1\bigr)^2
\]
is a perfect square.

\begin{problem1}
Are there constants $S$ and $M$ such that, for every $k$, some choice has no colliding ordered pair?
\end{problem1}

What I know.
\begin{itemize}
\item Collisions have a simple form. Write $Y=|2(j_v-i_u)-4s+1|$ and $q=|c(v)-c(u)|$. The pair
collides if and only if $Y=gef$ and $q=g(f^2-e^2)/(2S)$ for odd integers $g,e,f$ with $e<f$
coprime. For $S\ge4$ a power of 2 this needs $f\equiv\pm e\pmod S$.
\item For $S=32$ and $M=3$ a choice with no collision exists for $s=1024$. I found it by
computer and checked all $2{,}169{,}330$ candidate pairs. For $S=32$ and $M=2$ no choice exists
for $s=256$, by an odd cycle of 51 pairs.
\item The number of candidate partners of a cell grows. I sampled 2000 cells for each
$s=2^9,\dots,2^{16}$. The mean is about $(21.5/S)\log_2 s$ plus a constant, for $S=32$ and
$S=64$. Hence no argument through bounded degree can answer the problem.
\item The arithmetic alone does not help. Every finite colouring of the odd positive integers has
distinct $u,v$ of one colour with $u^2+4(v-u)^2$ a square. I proved this by adapting the work of
Frantzikinakis, Klurman and Moreira on partition regularity of Pythagorean pairs. A positive
answer must use the Hilbert order.
\end{itemize}

A negative answer would suggest that my method cannot reach $m_h=O(h^2)$. With $S$ of order
$\log s$ the degree stays bounded, and the method then aims at $h^2\log h$. I would also like
to know whether $m_h/h^2$ is unbounded. For large $h$ I have no lower bound for $m_h$ above $q_h$.

\section*{Problem 2. All-double windows}

Let $h\ge4$ and let $(A,B)$ be a partition of the rows $\{1,\dots,h\}$ with
$|A|=\lceil h/2\rceil$ and $|B|=\lfloor h/2\rfloor$. A \emph{window} of width $k$ is a set $P$ of
$2k$ points of the $h\times k$ grid with 3 properties. First, every column has 2 points of $P$,
one in a row of $A$ and one in a row of $B$. Second, the $k$ row pairs are distinct. Third, every
tilted lattice rectangle with its 4 corners in the $h\times k$ grid has at most 2 corners in $P$.
Let $W(h)$ be the largest width of a window, over all such partitions.

\begin{problem2}
Prove that $W(h)\le h$ for every $h\ge4$.
\end{problem2}

Why it matters. In a blackout of size $m+q_h$, every run of consecutive columns with 2 black
points is a window. Hence
\[
 m_h\ge q_h+\left\lfloor\frac{q_h-1}{W(h)}\right\rfloor .
\]
So $W(h)\le h$ gives $m_h\ge q_h+h/4-O(1)$, and even $W(h)<q_h$ gives $m_h>q_h$. With $W(h)=h$
for $h\le12$ the bound gives $m_h>q_h$ for $5\le h\le12$. For large $h$ I cannot prove $m_h>q_h$.

What I know.
\begin{itemize}
\item $W(h)\ge h$ for every $h\ge4$. For even $h$ take the columns
$\{1,2\},\{2,3\},\dots,\{h-1,h\},\{1,h\}$ with the parity partition. For odd $h$ take
$\{1,h\},\{1,2\},\{2,3\},\dots,\{h-2,h-1\},\{1,h-1\}$.
\item $W(h)=h$ for $4\le h\le12$, by SAT search. Each upper bound has an LRAT proof that 2 checkers accept. One program wrote the clauses, and a second program has confirmed the values only for $h\le10$.
\item A fractional argument cannot work. Put weight $1/|A|$ on every cell in a row of $A$ and $1/|B|$ on every cell in a row of $B$. This meets every condition at every width. So a proof must use integrality.
\item A window can have more columns than rows in use. For $h=12$ there is a window of width 9 on 6 rows.
\item Suppose a window contains the graph of a walk $f$ with $|f(x+1)-f(x)|=1$ across all its
columns. Then its width is at most $h$.
\item Rectangles with sides of slope $\pm1$ are not enough. A point set on 2 adjacent lines of
slope 1 has no tilted right angle, and sets that meet only the slope $\pm1$ condition can be
wider than $h$ for $h=5,6,9$.
\item A right angle whose 2 legs point to the same side has its fourth corner at most
$\lfloor(h-1)/2\rfloor$ columns beyond the farther leg. So the problem reduces, with a loss of
$h$ columns, to sets with no such right angle at all.
\end{itemize}

The windows of width $h$ for $h\le10$ are not all close to 2 adjacent lines. They have many
points in the top and bottom rows.

\section*{Files}

Notes, programs and certificates are at
\url{https://github.com/apemm/rectangle-blackouts}. I would be glad to hear of any related
result, in particular on point sets with no right angle and on colourings that avoid Pythagorean
pairs along a space-filling curve. The closest work I know is that of Shkredov and Solymosi,
\emph{Tilted corners in integer grids}, Integers 21A (2021), Paper A20.

\end{document}
